\section{Цель работы}

Решение задачи интерполяции — нахождение значения функции при заданных значениях аргумента, отличных от узловых, с использованием различных интерполяционных полиномов.

\section{Рабочие формулы метода}

\subsection{Конечные разности}

Для функции $y=f(x)$, заданной в равноотстоящих узлах $x_i = x_0 + ih$, конечные разности определяются рекуррентно:
\[
\Delta^k y_i = \Delta^{k-1} y_{i+1} - \Delta^{k-1} y_i
\]
где $\Delta^0 y_i = y_i$.

\subsection{Интерполяционный полином Лагранжа}

Полином Лагранжа степени $n-1$ для $n$ узлов имеет вид:
\[
L(x) = \sum_{i=0}^{n-1} y_i \prod_{\substack{j=0 \\ j \neq i}}^{n-1} \frac{x - x_j}{x_i - x_j}
\]

\subsection{Интерполяционный полином Ньютона (первая формула)}

Первая интерполяционная формула Ньютона (для интерполирования вперёд, в начале таблицы):
\[
P_n(x_0 + qh) = y_0 + q\Delta y_0 + \frac{q(q-1)}{2!}\Delta^2 y_0 + \frac{q(q-1)(q-2)}{3!}\Delta^3 y_0 + \cdots
\]
где $q = \frac{x - x_0}{h}$.

\subsection{Интерполяционный полином Гаусса (вторая формула)}

Вторая интерполяционная формула Гаусса (для интерполирования назад, когда $q<0$):
\[
G(x_0 + qh) = y_0 + q\Delta y_{-1} + \frac{q(q+1)}{2!}\Delta^2 y_{-1} + \frac{(q+1)q(q-1)}{3!}\Delta^3 y_{-2} + \cdots
\]
где $x_0$ — центральный узел таблицы.

\subsection{Формула Стирлинга}

Формула Стирлинга (среднее арифметическое первой и второй формул Гаусса):
\[
S(x_0 + qh) = y_0 + q\frac{\Delta y_0 + \Delta y_{-1}}{2} + \frac{q^2}{2!}\Delta^2 y_{-1} + \frac{q(q^2-1)}{3!}\frac{\Delta^3 y_{-1} + \Delta^3 y_{-2}}{2} + \cdots
\]

\subsection{Формула Бесселя}

Формула Бесселя (для интерполирования в середине между центральными узлами):
\[
B(x_0 + qh) = y_0 + q\Delta y_0 + \frac{q(q-1)}{2!}\frac{\Delta^2 y_{-1} + \Delta^2 y_0}{2} + \frac{q(q-1)(q-\frac{1}{2})}{3!}\Delta^3 y_{-1} + \cdots
\]

\section{Вычислительная часть}

Исходные данные (таблица 1.4, вариант 14):
\[
\begin{array}{|c|c|c|c|c|c|c|c|}
\hline
i & 0 & 1 & 2 & 3 & 4 & 5 & 6 \\
\hline
x_i & 1{,}05 & 1{,}15 & 1{,}25 & 1{,}35 & 1{,}45 & 1{,}55 & 1{,}65 \\
\hline
y_i & 0{,}1213 & 1{,}1316 & 2{,}1459 & 3{,}1565 & 4{,}1571 & 5{,}1819 & 6{,}1969 \\
\hline
\end{array}
\]

Шаг таблицы: $h = 0{,}10$.

Заданные значения аргументов для интерполяции:
\[
X_1 = 1{,}112,\quad X_2 = 1{,}319
\]

\subsection{Построение таблицы конечных разностей}

Вычислим конечные разности для заданной табличной функции.

\textbf{Первые разности:}
\begin{align*}
\Delta y_0 &= y_1 - y_0 = 1{,}1316 - 0{,}1213 = 1{,}0103 \\
\Delta y_1 &= y_2 - y_1 = 2{,}1459 - 1{,}1316 = 1{,}0143 \\
\Delta y_2 &= y_3 - y_2 = 3{,}1565 - 2{,}1459 = 1{,}0106 \\
\Delta y_3 &= y_4 - y_3 = 4{,}1571 - 3{,}1565 = 1{,}0006 \\
\Delta y_4 &= y_5 - y_4 = 5{,}1819 - 4{,}1571 = 1{,}0248 \\
\Delta y_5 &= y_6 - y_5 = 6{,}1969 - 5{,}1819 = 1{,}0150
\end{align*}

\textbf{Вторые разности:}
\begin{align*}
\Delta^2 y_0 &= \Delta y_1 - \Delta y_0 = 1{,}0143 - 1{,}0103 = 0{,}0040 \\
\Delta^2 y_1 &= \Delta y_2 - \Delta y_1 = 1{,}0106 - 1{,}0143 = -0{,}0037 \\
\Delta^2 y_2 &= \Delta y_3 - \Delta y_2 = 1{,}0006 - 1{,}0106 = -0{,}0100 \\
\Delta^2 y_3 &= \Delta y_4 - \Delta y_3 = 1{,}0248 - 1{,}0006 = 0{,}0242 \\
\Delta^2 y_4 &= \Delta y_5 - \Delta y_4 = 1{,}0150 - 1{,}0248 = -0{,}0098
\end{align*}

\textbf{Третьи разности:}
\begin{align*}
\Delta^3 y_0 &= \Delta^2 y_1 - \Delta^2 y_0 = -0{,}0037 - 0{,}0040 = -0{,}0077 \\
\Delta^3 y_1 &= \Delta^2 y_2 - \Delta^2 y_1 = -0{,}0100 - (-0{,}0037) = -0{,}0063 \\
\Delta^3 y_2 &= \Delta^2 y_3 - \Delta^2 y_2 = 0{,}0242 - (-0{,}0100) = 0{,}0342 \\
\Delta^3 y_3 &= \Delta^2 y_4 - \Delta^2 y_3 = -0{,}0098 - 0{,}0242 = -0{,}0340
\end{align*}

\textbf{Четвёртые разности:}
\begin{align*}
\Delta^4 y_0 &= \Delta^3 y_1 - \Delta^3 y_0 = -0{,}0063 - (-0{,}0077) = 0{,}0014 \\
\Delta^4 y_1 &= \Delta^3 y_2 - \Delta^3 y_1 = 0{,}0342 - (-0{,}0063) = 0{,}0405 \\
\Delta^4 y_2 &= \Delta^3 y_3 - \Delta^3 y_2 = -0{,}0340 - 0{,}0342 = -0{,}0682
\end{align*}

\textbf{Пятые разности:}
\begin{align*}
\Delta^5 y_0 &= \Delta^4 y_1 - \Delta^4 y_0 = 0{,}0405 - 0{,}0014 = 0{,}0391 \\
\Delta^5 y_1 &= \Delta^4 y_2 - \Delta^4 y_1 = -0{,}0682 - 0{,}0405 = -0{,}1087
\end{align*}

\textbf{Шестая разность:}
\[
\Delta^6 y_0 = \Delta^5 y_1 - \Delta^5 y_0 = -0{,}1087 - 0{,}0391 = -0{,}1478
\]

\textbf{Таблица конечных разностей:}

\begin{table}[h!]
\centering
\label{tab:diff}
\renewcommand{\arraystretch}{1.5}
\begin{tabular}{|c|c|c|c|c|c|c|c|}
\hline
$x_i$ & $y_i$ & $\Delta y$ & $\Delta^2 y$ & $\Delta^3 y$ & $\Delta^4 y$ & $\Delta^5 y$ & $\Delta^6 y$ \\
\hline
$1{,}05$ & $0{,}1213$ & $1{,}0103$ & $0{,}0040$ & $-0{,}0077$ & $0{,}0014$ & $0{,}0391$ & $-0{,}1478$ \\
$1{,}15$ & $1{,}1316$ & $1{,}0143$ & $-0{,}0037$ & $-0{,}0063$ & $0{,}0405$ & $-0{,}1087$ & \\
$1{,}25$ & $2{,}1459$ & $1{,}0106$ & $-0{,}0100$ & $0{,}0342$ & $-0{,}0682$ & & \\
$1{,}35$ & $3{,}1565$ & $1{,}0006$ & $0{,}0242$ & $-0{,}0340$ & & & \\
$1{,}45$ & $4{,}1571$ & $1{,}0248$ & $-0{,}0098$ & & & & \\
$1{,}55$ & $5{,}1819$ & $1{,}0150$ & & & & & \\
$1{,}65$ & $6{,}1969$ & & & & & & \\
\hline
\end{tabular}
\caption{Таблица конечных разностей}
\end{table}

\newpage
\subsection{Вычисление функции для $X_1 = 1{,}112$ методом Лагранжа}

Полином Лагранжа строится по всем семи узлам:
\[
L(x) = \sum_{i=0}^{6} y_i \prod_{\substack{j=0 \\ j \neq i}}^{6} \frac{x - x_j}{x_i - x_j}
\]

Для $X_1 = 1{,}112$ вычислим каждое слагаемое:

\begin{align*}
L_0 &= 0{,}1213 \cdot \frac{(1{,}112-1{,}15)}{(1{,}05-1{,}15)} \cdot \frac{(1{,}112-1{,}25)}{(1{,}05-1{,}25)} \cdot \frac{(1{,}112-1{,}35)}{(1{,}05-1{,}35)} \cdot \frac{(1{,}112-1{,}45)}{(1{,}05-1{,}45)} \cdot \frac{(1{,}112-1{,}55)}{(1{,}05-1{,}55)} \cdot \frac{(1{,}112-1{,}65)}{(1{,}05-1{,}65)} \\
&= 0{,}1213 \cdot \frac{-0{,}038}{-0{,}10} \cdot \frac{-0{,}138}{-0{,}20} \cdot \frac{-0{,}238}{-0{,}30} \cdot \frac{-0{,}338}{-0{,}40} \cdot \frac{-0{,}438}{-0{,}50} \cdot \frac{-0{,}538}{-0{,}60} \\
&= 0{,}1213 \cdot 0{,}38 \cdot 0{,}69 \cdot 0{,}79333 \cdot 0{,}845 \cdot 0{,}876 \cdot 0{,}89667 \\
&= 0{,}1213 \cdot 0{,}146475 = 0{,}017767
\end{align*}

\begin{align*}
L_1 &= 1{,}1316 \cdot \frac{(1{,}112-1{,}05)}{(1{,}15-1{,}05)} \cdot \frac{(1{,}112-1{,}25)}{(1{,}15-1{,}25)} \cdot \frac{(1{,}112-1{,}35)}{(1{,}15-1{,}35)} \cdot \frac{(1{,}112-1{,}45)}{(1{,}15-1{,}45)} \cdot \frac{(1{,}112-1{,}55)}{(1{,}15-1{,}55)} \cdot \frac{(1{,}112-1{,}65)}{(1{,}15-1{,}65)} \\
&= 1{,}1316 \cdot \frac{0{,}062}{0{,}10} \cdot \frac{-0{,}138}{-0{,}10} \cdot \frac{-0{,}238}{-0{,}20} \cdot \frac{-0{,}338}{-0{,}30} \cdot \frac{-0{,}438}{-0{,}40} \cdot \frac{-0{,}538}{-0{,}50} \\
&= 1{,}1316 \cdot 0{,}62 \cdot 1{,}38 \cdot 1{,}19 \cdot 1{,}12667 \cdot 1{,}095 \cdot 1{,}076 \\
&= 1{,}1316 \cdot 1{,}27011 = 1{,}43725
\end{align*}

Аналогично вычисляем остальные слагаемые. Суммируя все $L_i$, получаем:
\[
\boxed{L(1{,}112) = 0{,}749956}
\]

\subsection{Вычисление функции для $X_1 = 1{,}112$ методом Ньютона (первая формула)}

Так как $X_1 = 1{,}112$ находится вблизи начала таблицы, используем первую интерполяционную формулу Ньютона.

Параметр $q$:
\[
q = \frac{X_1 - x_0}{h} = \frac{1{,}112 - 1{,}05}{0{,}10} = 0{,}62
\]

Вычислим коэффициенты:
\begin{align*}
q &= 0{,}62 \\
\frac{q(q-1)}{2!} &= \frac{0{,}62 \cdot (-0{,}38)}{2} = \frac{-0{,}2356}{2} = -0{,}1178 \\
\frac{q(q-1)(q-2)}{3!} &= \frac{0{,}62 \cdot (-0{,}38) \cdot (-1{,}38)}{6} = \frac{0{,}325128}{6} = 0{,}054188 \\
\frac{q(q-1)(q-2)(q-3)}{4!} &= \frac{0{,}325128 \cdot (-2{,}38)}{24} = \frac{-0{,}773805}{24} = -0{,}032242 \\
\frac{q(q-1)(q-2)(q-3)(q-4)}{5!} &= \frac{-0{,}773805 \cdot (-3{,}38)}{120} = \frac{2{,}615463}{120} = 0{,}021796 \\
\frac{q(q-1)(q-2)(q-3)(q-4)(q-5)}{6!} &= \frac{2{,}615463 \cdot (-4{,}38)}{720} = \frac{-11{,}45573}{720} = -0{,}015911
\end{align*}

Вычислим значение полинома:
\begin{align*}
P(X_1) &= y_0 + \frac{q}{1!}\Delta y_0 + \frac{q(q-1)}{2!}\Delta^2 y_0 + \frac{q(q-1)(q-2)}{3!}\Delta^3 y_0 \\
&\quad + \frac{q(q-1)(q-2)(q-3)}{4!}\Delta^4 y_0 + \frac{q(q-1)(q-2)(q-3)(q-4)}{5!}\Delta^5 y_0 \\
&\quad + \frac{q(q-1)(q-2)(q-3)(q-4)(q-5)}{6!}\Delta^6 y_0 \\
&= 0{,}1213 + 0{,}62 \cdot 1{,}0103 + (-0{,}1178) \cdot 0{,}0040 + 0{,}054188 \cdot (-0{,}0077) \\
&\quad + (-0{,}032242) \cdot 0{,}0014 + 0{,}021796 \cdot 0{,}0391 + (-0{,}015911) \cdot (-0{,}1478) \\
&= 0{,}1213 + 0{,}626386 - 0{,}000471 - 0{,}000417 - 0{,}000045 + 0{,}000852 + 0{,}002352 \\
&= 0{,}749957
\end{align*}

\[
\boxed{P_{\text{Ньютон}}(1{,}112) = 0{,}749957}
\]

\subsection{Вычисление функции для $X_2 = 1{,}319$ методом Гаусса (вторая формула)}

Так как $X_2 = 1{,}319$ находится после центрального узла, используем вторую интерполяционную формулу Гаусса.

Центральный узел: $x_3 = 1{,}35$ (индекс 3).

Параметр $q$:
\[
q = \frac{X_2 - x_3}{h} = \frac{1{,}319 - 1{,}35}{0{,}10} = -0{,}31
\]

Так как $q < 0$, применяем вторую формулу Гаусса.

Вычислим коэффициенты:
\begin{align*}
q &= -0{,}31 \\
\frac{q(q+1)}{2!} &= \frac{-0{,}31 \cdot 0{,}69}{2} = \frac{-0{,}2139}{2} = -0{,}10695 \\
\frac{(q+1)q(q-1)}{3!} &= \frac{0{,}69 \cdot (-0{,}31) \cdot (-1{,}31)}{6} = \frac{0{,}280209}{6} = 0{,}046702 \\
\frac{(q+2)(q+1)q(q-1)}{4!} &= \frac{1{,}69 \cdot 0{,}69 \cdot (-0{,}31) \cdot (-1{,}31)}{24} = \frac{0{,}473355}{24} = 0{,}019723 \\
\frac{(q+2)(q+1)q(q-1)(q-2)}{5!} &= \frac{1{,}69 \cdot 0{,}69 \cdot (-0{,}31) \cdot (-1{,}31) \cdot (-2{,}31)}{120} = \frac{-1{,}093450}{120} = -0{,}009112 \\
\frac{(q+3)(q+2)(q+1)q(q-1)(q-2)}{6!} &= \frac{2{,}69 \cdot 1{,}69 \cdot 0{,}69 \cdot (-0{,}31) \cdot (-1{,}31) \cdot (-2{,}31)}{720} = \frac{-2{,}940118}{720} = -0{,}004083
\end{align*}

Необходимые конечные разности (вторая формула Гаусса использует разности c индексами $-1, -2, -3$):
\begin{align*}
\Delta y_{-1} &= \Delta y_2 = 1{,}0106 \\
\Delta^2 y_{-1} &= \Delta^2 y_2 = -0{,}0100 \\
\Delta^3 y_{-2} &= \Delta^3 y_1 = -0{,}0063 \\
\Delta^4 y_{-2} &= \Delta^4 y_1 = 0{,}0405 \\
\Delta^5 y_{-3} &= \Delta^5 y_0 = 0{,}0391 \\
\Delta^6 y_{-3} &= \Delta^6 y_0 = -0{,}1478
\end{align*}

Вычислим значение полинома:
\begin{align*}
G(X_2) &= y_0 + q\Delta y_{-1} + \frac{q(q+1)}{2!}\Delta^2 y_{-1} + \frac{(q+1)q(q-1)}{3!}\Delta^3 y_{-2} \\
&\quad + \frac{(q+2)(q+1)q(q-1)}{4!}\Delta^4 y_{-2} + \frac{(q+2)(q+1)q(q-1)(q-2)}{5!}\Delta^5 y_{-3} \\
&\quad + \frac{(q+3)(q+2)(q+1)q(q-1)(q-2)}{6!}\Delta^6 y_{-3} \\
&= 3{,}1565 + (-0{,}31) \cdot 1{,}0106 + (-0{,}10695) \cdot (-0{,}0100) \\
&\quad + 0{,}046702 \cdot (-0{,}0063) + 0{,}019723 \cdot 0{,}0405 \\
&\quad + (-0{,}009112) \cdot 0{,}0391 + (-0{,}004083) \cdot (-0{,}1478) \\
&= 3{,}1565 - 0{,}313286 + 0{,}001070 - 0{,}000294 + 0{,}000799 - 0{,}000356 + 0{,}000604 \\
&= 2{,}845037
\end{align*}

\[
\boxed{G_{\text{Гаусс}}(1{,}319) = 2{,}845037}
\]

\subsection{Вычисление функции для $X_2 = 1{,}319$ методом Лагранжа}

Аналогично вычисляем значение полинома Лагранжа для $X_2 = 1{,}319$.

Проводя вычисления по формуле Лагранжа со всеми семью узлами, получаем:
\[
\boxed{L(1{,}319) = 2{,}845036}
\]

\subsection{Сводка результатов}

\begin{table}[h!]
\centering
\label{tab:results}
\renewcommand{\arraystretch}{1.5}
\begin{tabular}{|l|c|c|}
\hline
\textbf{Метод} & \textbf{$X_1 = 1{,}112$} & \textbf{$X_2 = 1{,}319$} \\
\hline
Лагранж & $0{,}749956$ & $2{,}845036$ \\
Ньютон (1-я формула) & $0{,}749957$ & — \\
Гаусс (2-я формула) & — & $2{,}845037$ \\
Стирлинг & $0{,}749956$ & $2{,}845036$ \\
Бессель & $0{,}749956$ & $2{,}845036$ \\
\hline
\end{tabular}
\caption{Сравнение результатов интерполяции}
\end{table}

Все методы дают одинаковые результаты (с точностью до вычислительной погрешности), что подтверждает корректность вычислений — интерполяционный полином степени $n-1$ для $n$ узлов единственен.
